% ===========================================================================
% Hardware Realization -- REFRAMED (drop-in for \section{Hardware Realization}
% \label{sec:Proposed}). Changes vs. the prior draft, all surgical:
%   1. Fixed the inverted-thesis sentence ("we do claim" -> "we do NOT claim").
%   2. GF4 is modeled ONLY at fp16 (gf4_real, Table~\ref{tab:gf4-energy}); int4
%      appears solely as the NVFP4/E2M1 baseline, never on GF4. The decode is
%      isolated multiply-free (MAC removed from both engines), so no number
%      charges GF4 a 4-bit multiply (int4-as-GF4 tables discarded).
%   3. Named the dense-fp16 CUDA baseline (Table~\ref{tab:cuda-speedup}) beside
%      the fake-quant naive baseline.
%   4. Folded the CNN rows of Tables~\ref{tab:lut},~\ref{tab:gf4-energy} in.
%   5. LUTxLUT figure reconciled to the synth log (9376 vs 2180 um^2; joint
%      table 8877 um^2) and framed as the validated dead-end that confirms
%      decode-then-multiply.
%   6. Typo "frame the of this section" -> "frame the reading of this section".
% ===========================================================================

% To improve the hardware compatibility of Gaussian-aware 4-bit quantization, we use the codebook
% \(\mathcal{C}=\{0,1,3,5,6,8,11,16\}/16\), which approximates the distribution-aware levels of NormalFloat while constraining all 
% reconstruction values to a fixed \(1/16\) lattice. Although prior work has explored distribution-aware 4-bit codebooks such as NF4, 
% optimized variants such as AF4 and BOF4, and hardware-efficient nonuniform representations such as APoT and SBVR 
% \cite{dettmers2023qlora,yoshida2023af4,blumenberg2026bof4,li2020apot,bang2025sbvr}, we found no prior work using this specific Gaussian-derived 
% fixed-point codebook. This distinguishes our approach from methods that employ arbitrary floating-point lookup tables, such as any4 
% \cite{elhoushi2025any4}, and motivates its potential for efficient implementation across GPUs without relying on Blackwell-specific numerical formats.


% @article{dettmers2023qlora,
%   title   = {QLoRA: Efficient Finetuning of Quantized LLMs},
%   author  = {Dettmers, Tim and Pagnoni, Artidoro and Holtzman, Ari and Zettlemoyer, Luke},
%   journal = {Advances in Neural Information Processing Systems},
%   volume  = {36},
%   year    = {2023}
% }

% @article{yoshida2023af4,
%   title   = {Affine4: A New 4-Bit Quantization Method for Large Language Models},
%   author  = {Yoshida, Gakuto},
%   journal = {arXiv preprint arXiv:2306.06965},
%   year    = {2023}
% }

% @inproceedings{blumenberg2026bof4,
%   title     = {Improving Block-Wise LLM Quantization by 4-bit Block-Wise Optimal Float (BOF4): Analysis and Variations},
%   author    = {Blumenberg, Patrick and Graave, Thomas and Fingscheidt, Tim},
%   booktitle = {International Conference on Learning Representations},
%   year      = {2026}
% }

% @inproceedings{li2020apot,
%   title     = {Additive Powers-of-Two Quantization: An Efficient Non-uniform Discretization for Neural Networks},
%   author    = {Li, Yuhang and Dong, Xin and Wang, Wei},
%   booktitle = {International Conference on Learning Representations},
%   year      = {2020}
% }

% @article{bang2025sbvr,
%   title   = {SBVR: Summation of BitVector Representation for Efficient LLM Quantization},
%   author  = {Bang, Wonjun and Park, Jongseok and Yu, Hongseung and Bin, Kyungmin and Lee, Kyunghan},
%   journal = {arXiv preprint arXiv:2509.18172},
%   year    = {2025}
% }

% @inproceedings{elhoushi2025any4,
%   title     = {any4: Learned 4-bit Numeric Representation for LLMs},
%   author    = {Elhoushi, Mostafa and Johnson, Jeff},
%   booktitle = {Proceedings of the 42nd International Conference on Machine Learning},
%   volume    = {267},
%   pages     = {15215--15231},
%   year      = {2025},
%   publisher = {PMLR}
% }

\section{Hardware Realization}\label{sec:Proposed}

We propose and simulate an inexpensive dedicated hardware unit to support both
codebook flexibility and outlier robustness. This includes modeling a
weight-stationary FP4 systolic array in Timeloop and Accelergy
\cite{parashar2019timeloop,wu2019accelergy}, evaluating and confirming these
models with an RTL implementation in Verilog, and then evaluating two things:
(1) the cost of a fully programmable decode table, and (2) the cost of serving
outlier layers with multiple passes compared to a dedicated FP16 unit.

Three claims are being made in this section, and they are evaluated separately
because they rest on different evidence. The first is that a programmable
codebook, the mechanism that lets the same hardware decode either NVFP4 or GF4,
costs almost nothing if it is placed in the right spot in the datapath.
Section \ref{subsec:acc} and Section \ref{subsec:progcode} establish this with a
Timeloop/Accelergy component model. The second is that outlier layers can be
served by summing multiple GF4 passes into one operand and multiplying once
through a modestly widened multiplier, rather than adding a full second,
16-bit-wide arithmetic datapath, covered in Section \ref{subsec:multipass}.
Section \ref{sub:rtl} then checks both against real evidence beyond the component
model with a synthesized RTL prototype and measured CUDA kernels --- and, in
doing so, finds that the two ways of executing GF4's arbitrary magnitudes (a
product-table datapath and GPU decode-to-fp16) are each unfruitful as a
hardware-native format. That negative result motivates the third and central
claim: the codebook itself can be constrained to a fixed-point lattice and
\emph{optimized} within that constraint, yielding a format that is cheap to
decode by construction, universal across models, and iso-accurate with the
hand-tuned Gaussian codebook (Section \ref{subsec:codebook-design}).

The mechanism for this proposed hardware is as follows. A conventional FP4
accelerator decodes each 4-bit code by looking it up in a small, fixed table of
magnitudes (E2M1's eight values) before feeding it to a multiplier. We keep that
exact decode-then-multiply structure and change exactly one thing.
Specifically, the table's contents are no longer hardwired, but held in a small
writable register file, so the same array can decode either NVFP4's magnitudes
or GF4's. Section \ref{subsec:acc} quantifies what that one change costs in a
realistic accelerator model. Section \ref{subsec:progcode} shows where in the
pipeline to place the decode so that cost is nearly zero; Section
\ref{subsec:multipass} covers the separate mechanism that handles outlier
layers. Section \ref{sub:rtl} confirms both mechanisms in real RTL and on real
GPU hardware. Note that results for the aforementioned sections characterize the
single-pass (T=1) engine, and the widened multiplier residual GF4 requires for
T$>$1 is not separately synthesized or measured here. A note of realism helps to
frame the reading of this section. GF4's magnitude levels are not four-bit
values, so no genuine four-bit GF4 multiplier exists. The 4-bit code indexes an
eight-entry codebook of FP4's cardinality, whose magnitudes are held at higher
precision and feed an ordinary ($\approx$16-bit) multiplier. The accelerator's
contribution is therefore not a narrower multiplier but a nearly free,
programmable \emph{decode}. The component model prices GF4's arithmetic
faithfully --- an fp16 multiply on the decoded $\approx$16-bit magnitude --- and
reports GF4's total cost, decode plus fp16 multiply, directly against the NVFP4
baseline (Table~\ref{tab:gf4-energy}); to isolate the decode alone it removes the
multiplier from both engines and prices the codebook against the memory system
(Table~\ref{tab:lut}), a number that does not depend on how arithmetic is
modeled. GF4 is never modeled on a 4-bit multiplier: its levels are not 4-bit
values, so that would misstate its cost. Because the decode-then-multiply
structure already runs on every deployed GPU, the immediately usable realization
is the CUDA kernels of Section~\ref{sec:sw-support}, not a bespoke chip.


\subsection{Accelerator Model}\label{subsec:acc}

In Timeloop with Accelergy component energy/area (45 nm, Aladdin/CACTI) we model
a weight-stationary 16$\times$16 (256-PE) systolic array (DRAM $\rightarrow$
global buffer $\rightarrow$ per-PE register files $\rightarrow$ MAC). The GF4 PE
--- the programmable-decode datapath proposed above --- pairs a 16-entry codebook
register file with an fp16 multiplier: because GF4's levels are not four-bit
values, the decoded operand is a $\approx$16-bit magnitude and the multiply is
genuinely fp16, which is exactly how we model it (engine \texttt{gf4\_real}). The NVFP4 baseline is the identical array with the
codebook removed and a 4-bit fixed-point MAC --- the faithful low-energy proxy
for a native E2M1 unit (Accelergy cannot represent E2M1; see below). Modeling
GF4 at fp16 and NVFP4 at int4 is the honest pairing: it charges each format for
the multiplier it actually uses, and never prices GF4 on a 4-bit MAC it does not.
The decode is then isolated \emph{multiply-free} (Table~\ref{tab:lut}), removing
the MAC from both engines so the codebook's cost stands on its own.

Weight-stationary dataflow keeps decoded weights resident and reused across
activation columns, which is why decode-at-ingress is natural here and why the
placement is nearly free. It is also the canonical GEMM dataflow (Eyeriss, TPU),
making the baseline uncontroversial. 256 PEs is the smallest array where a
realistically provisioned global buffer dominates area. Smaller arrays are
governed by per-PE and array-edge overheads Accelergy models only coarsely,
whereas larger arrays require proportionally larger buffers, leaving the
codebook's area share unchanged or smaller. All claims are ratios (codebook
overhead per PE, energy/token), so they are scale-invariant. We stay at the
45 nm node so every number traces to a validated CACTI/Aladdin library rather
than an advanced-node extrapolation. The relative ordering is node-robust.
Two engines are simulated directly in Timeloop/Accelergy: \texttt{gf4\_real} ---
our proposed datapath, a programmable codebook decode feeding an fp16 multiply ---
and a native NVFP4 baseline (table-free, int4 MAC; identical 2.394 mm$^2$ array
area). These differ in \emph{both} the decode (table vs.\ none) \emph{and} the
multiplier (fp16 vs.\ int4), so their raw difference is the realistic total cost
of Table~\ref{tab:gf4-energy}, decode plus wider multiply. To isolate the codebook
alone we remove the multiplier from both engines --- the multiply-free figure of
Table~\ref{tab:lut}. This is also precisely the picture \emph{on the proposed
datapath itself}: there a single fp16 multiplier serves either format and only the
writable table's contents change, so switching a layer from NVFP4 to GF4 costs
nothing but that table, and the multiply-free number \emph{is} that switching cost.
The codebook's own energy is read from the simulated decode component, and the
decode-at-ingress placement is priced analytically from these same runs' access
counts and cross-checked against the RTL decode synthesis
(Table~\ref{tab:node-scaling}).

\paragraph{Faithful multiply cost.} With GF4 modeled at fp16 throughout, the
wider multiply is measured rather than assumed. Holding the decode placement
fixed, the fp16-multiply premium is $11.3$--$20.0\%$ on the compute-bound language models
and $0$--$3.3\%$ on the memory-bound CNNs, whose energy is dominated by memory
traffic; the realistic total GF4 overhead (decode $+$ fp16 multiply) is
$32.6$--$42.7\%$ (LLMs) and $3.4$--$20.8\%$ (CNNs), reported in
Table~\ref{tab:gf4-energy}. One modeling caveat must be stated plainly: Accelergy's
Aladdin models cannot represent an E2M1 (2-exponent, 1-mantissa) multiplier ---
the \texttt{fpmac} primitive floors near fp8, so a literal E2M1 unit reports
\emph{higher} energy than a 4-bit integer MAC, which is physically backwards. We
therefore keep the NVFP4 baseline modeled as a 4-bit fixed-point MAC, the
faithful low-energy proxy for a native E2M1 unit; a true E2M1 unit is no cheaper,
so the premiums above are, if anything, conservative.\footnote{A hardware run
with a custom-characterized E2M1 multiplier would remove this proxy and tighten
the multiply-cost figures; the decode-isolation and speedup conclusions are
unaffected by it.}


\subsection{Programmable-Codebook Cost}\label{subsec:progcode}

The energy model prices a conservative 16-entry (full 4-bit) table, though GF4
needs only 8 magnitude entries plus a sign bit. The MAC and a 16-entry table
read each cost $\approx$ 1 pJ at 45 nm (Aladdin/CACTI, c.f Chapter
\ref{chap:refsysmod}), so the energy cost of the table is set by how often it is
read, which in turn is set by where in the datapath it fires.

Decode-per-MAC fires on every multiply and reads both operands' tables, so its
cost is priced at 2 $\times$ Computes table reads, independent of reuse. We report
this decode cost \emph{multiply-free} --- the codebook's energy against the
memory system, with the multiplier removed from both engines --- so it does not
depend on how arithmetic is modeled and never charges GF4 a 4-bit multiply. On
that basis the decode costs 15.71\% (OPT-125M), 25.62\% (OPT-1.3B), and 23.75\%
(Llama-2-7B) of the memory system (Table \ref{tab:lut}). This is not a monotonic
function of model or hidden-dimension size. Based upon our measurements
Llama-2-7B's overhead is slightly below OPT-1.3B's despite a larger hidden
dimension. What it does track is the share of total energy spent on compute
versus memory traffic. OPT-125M is the most memory-bound of the three, so its
memory system is the largest fraction of the total and the codebook's share of
it is smallest. In contrast, OPT-1.3B and Llama-2-7B are both more
compute-energy-dominated and land in a similar, higher band. We don't have (and
don't claim) a clean scaling law linking the overhead to a single architectural
parameter across all three points.

Compared to Decode-per-MAC, Decode-at-ingress instead fires once per operand
loaded from DRAM into the global buffer. Because a weight-stationary array reuses
each loaded operand across many MACs, the table is read only \verb|DRAM_accesses|
times rather than 2 $\times$ Computes times. We price this placement analytically
rather than simulating it directly as this is a methodological necessity, not a
choice. The decode-per-MAC attaches the LUT to a compute compound (MAC + table),
which Accelergy models correctly. Conversely, decode-at-ingress attaches the LUT
to a storage compound (the global-buffer SRAM), and wrapping a storage element in
an Accelergy compound triggers a known estimation bug that collapses its
area/energy (an initial direct simulation attempt returned nonsense figures such
as -80\% area). Timeloop's event-counting is unaffected by this bug, so we
instead take the DRAM access counts from the same, reliable NVFP4 Timeloop run
and price the decode cost by hand at \verb|E_LUT| = 1 pJ per 16-entry read,
grounded to a single regfile read at 45 nm. This analytical method
(\verb|accesses| $\times$ \verb|E_LUT|) reproduces the simulated codebook-decode
component energy to the decimal, and is independently corroborated by the RTL
decode synthesis (Table~\ref{tab:node-scaling}, $\approx$1.5\% of the PE). Priced
this way against the NVFP4 baseline, decode-at-ingress costs 1.08\%, 0.36\%, and
0.68\% for OPT-125M, OPT-1.3B, and Llama-2-7B respectively. These are decode-only
figures; GF4's fp16 multiply is measured separately and reported in
Table~\ref{tab:gf4-energy} (Section~\ref{subsec:acc}).

On that basis, decode-at-ingress costs 0.36--1.08\% energy and a flat 0.01\% area
across all three models tested. Reuse cuts the codebook cost sharply relative to
per-MAC decode --- the table fires once per DRAM load instead of twice per MAC, a
13--63$\times$ reduction across the three models. In other words, the codebook's
relative cost is small and, once placed at ingress, nearly vanishes. At-ingress trades a 2$\times$ larger
on-chip buffer for decoded operands (weights/activations are held decoded on-chip
rather than as 4-bit codes) but keeps DRAM traffic and the bandwidth-critical
path at 4-bit throughout.

The same pattern holds across convolutional networks
(Tables~\ref{tab:lut},~\ref{tab:gf4-energy}): the multiply-free decode ranges
from $3.5\%$ (LeNet-300-100) to $19\%$ (ResNet-56), and the fp16-multiply premium
stays $\le 3.3\%$ throughout, confirming the memory-bound regime in which GF4's
cost is decode rather than arithmetic.

Unlike per-tensor format selection (ANT, AdaptivFloat), which fixes the format
once chosen, GF4's table is re-decodable per access at the same amortized cost,
because decode-at-ingress ties cost to DRAM traffic, not to the number of formats
supported.


\subsection{Multi-Pass Precision}\label{subsec:multipass}

Residual GF4 generalizes to T passes, with each pass quantizing the residual of
previous passes. The T codes are summed in full precision before the multiply, so
the PE's multiplier is widened once to match the summed operand rather than being
replicated T times. Precision is dictated by pass count, traded against that
one-time multiplier widening rather than against an T-way LUT$\times$LUT product
table, which Section \ref{sub:rtl}'s RTL synthesis shows costs 4.3$\times$ more
area than a decode-then-multiply baseline. The residual of a GF4-quantized
Gaussian is itself near-Gaussian, so GF4 stays near-optimal on every subsequent
pass. The first residual pass cuts relative RMS error by roughly $8\times$
(10.16\% $\rightarrow$ 1.21\% on Gaussian input, c.f.\ Table~\ref{tab:nvfp4_noise}).


\subsection{RTL Verification}\label{sub:rtl}

The cost claims of Section \ref{subsec:progcode} combine a Timeloop/Accelergy
component model, measured directly for per-MAC and the NVFP4 baseline, with an
analytical extension of that same model's access counts for decode-at-ingress.
Both are validated pre-silicon methodology, but a model, not a circuit. This
section provides a synthesizeable RTL prototype to verify results.

We built and verified a staged RTL prototype of the datapath in synthesizable
Verilog, confirmed end-to-end on two independent simulators (Icarus Verilog and
Verilator) and cross-checked at each stage. Its most aggressive element is a
joint weight-code~$\times$~activation-code product table (``LUT$\times$LUT'') that
replaces decode-then-multiply outright. The LUT$\times$LUT is a single 64-entry
lookup ($8\times8$ magnitude products) plus a sign XOR, yielding the operand
product directly with no multiplier.

Correctness and area are separate claims, and the area results show that removal
of the multiplier does not save area. Synthesized against a fair
decode-then-multiply baseline (both operands decoded through an 8-entry
programmable table, then multiplied) with real RTL synthesis (Yosys/ABC,
Nangate45, not a component model), the 64-entry joint product table costs
9376\,\textmu m$^2$ against the baseline's 2180\,\textmu m$^2$. Area required for
LUT$\times$LUT is 4.3$\times$ larger, not smaller, at single-PE granularity. The
table's storage dominates (8877\,\textmu m$^2$ of the 9376 total), whereas a
9-bit signed multiplier is comparatively inexpensive in standard cells, and these
results match an independent finding from LUT Tensor Core~\cite{mo2025lut}:
conventional LUT-based mpGEMM designs lose their area advantage over MAC-based
designs once weight precision exceeds 2 bits, specifically because of table
precompute and storage overhead. Their fix is array-level table reuse and tiling
across many PEs, an optimization this single-PE RTL prototype does not implement.
We report this cost honestly rather than claim an unqualified win. At this scale,
LUT$\times$LUT's benefit is a genuinely multiplier-free datapath and the
falsifiable correctness guarantee verified above, not a demonstrated area
reduction. Crucially, this does not weaken GF4: GF4's datapath is
decode-then-multiply, which the same synthesis shows is the cheaper of the two,
so the LUT$\times$LUT result confirms rather than undermines the chosen design.

Due to the cost of replacing the multiplier with a product table, we retain the
multiplier to then examine the alternatives of various decode LUT forms. From RTL
synthesis, a fixed hardcoded GF4 decode incurred $15\,\mu\text{m}^2$ at 45\,nm,
roughly $1.5\%$ of a PE, and this metric remains constant across a node sweep of
sizes 130, 45, and 7\,nm ($1.61\%$, $1.46\%$, $1.57\%$; Table~\ref{tab:node-scaling}).
Writable decodes prove to be far more expensive, with a fully programmable
register file decode costing $599.56\,\mu\text{m}^2$, $40\times$ the fixed decode
size. When the writable table is folded into a complete PE, the marginal cost
falls to $+53\,\mu\text{m}^2$ per PE (a programmable-decode PE is
$1021\,\mu\text{m}^2$ versus the fixed-decode PE's $968\,\mu\text{m}^2$, a
$+5.5\%$ increase). Even at that reduced figure, the array has 256 PEs, so putting
a writable table in each one pays $+5.5\%$ 256 times over, for a table whose
contents are identical in every PE.

GF4 avoids the whole tax by committing to a hardwired decode, so it only pays the
$\approx$$1.5\%$ per-PE fixed cost and never instantiates a writable table inside
the array. In scenarios where the design genuinely needs to switch codebooks,
e.g. decode NVFP4 or GF4 on the same hardware, that single writable table should
be placed \emph{at ingress}. It is read once per operand as it loads from DRAM
into the buffer and then decoded once and reused across all of that operand's
MACs (Table~\ref{tab:lut}), rather than duplicated in every PE, in order to
amortize energy costs.

\paragraph{Measured GPU speedup.} Section~\ref{sec:sw-support} provides a
complementary, immediately deployable validation path: a suite of CUDA kernels
that run the same decode-then-multiply datapath on any modern NVIDIA GPU. Against
a dense all-fp16 GEMV --- native fp16 weights, the format deployed models
actually run --- the fused GF4 GEMV is $1.5$--$1.7\times$ faster in the
memory-bound regime (L2 flushed between calls, so weights stream from HBM as in a
real large-model layer; Table~\ref{tab:cuda-speedup}). This is the pure
weight-traffic win ($3.6\times$ fewer bytes) at an identical fp16 multiplier on
both sides, corroborating the component model's decode-and-traffic story on
production silicon.

Section~\ref{sec:sw-support} details these kernels; the important point for the
hardware argument is that the decode-then-multiply structure GF4 commits to is
not only cheap in the RTL and the component model, but already fast on GPUs in
the field.


\subsection{Hardware-Constrained Codebook Design}\label{subsec:codebook-design}

Section~\ref{sec:sw-support} closed on the one cost an arbitrary
distribution-matched codebook cannot shed: its decode is a lookup of irregular
magnitudes. The other obvious route to removing it is ruled out by the RTL of
Section~\ref{sub:rtl} --- a multiplication-free product table (LUT$\times$LUT) over
the codebook costs $4.3\times$ the PE area, because a table over arbitrary
magnitudes is large. A product table and GPU decode-to-fp16 both \emph{accept}
GF4's irrational levels and pay to work around them. We take the opposite move:
constrain the codebook itself to hardware-representable values and optimize
\emph{within} that constraint, so the format is cheap to decode by construction
rather than by emulation.

We therefore pose the codebook as the solution of a hardware-constrained,
distribution-aware optimization. For the operand distribution $p$ (the per-block
RMS-normalized magnitudes), the codebook $\mathcal{C}=\{\ell_0,\dots,\ell_{L-1}\}$
minimizing reconstruction distortion is
\begin{equation}
  \mathcal{C}^\star = \arg\min_{\mathcal{C}}\;
  \mathbb{E}_{w\sim p}\!\left[(w - Q_{\mathcal{C}}(w))^2\right],
  \label{eq:codebook-opt}
\end{equation}
with $Q_{\mathcal{C}}$ nearest-level quantization, subject to hardware
constraints on the levels: $\mathcal{C}\subset 2^{-m}\mathbb{Z}$ (representable on
an $m$-bit fractional grid), $|k_i|\le 2^{m}$ (bounded magnitude, top level
$=1$), and optionally $\operatorname{popcount}(k_i)\le r$ (few set bits, so decode
is shift-and-add). The magnitudes $\ell_i=k_i/2^m$ are integers on the grid; GF4
is the continuous ($m\!\to\!\infty$) Gaussian-quantile solution, and the deployed
format is one solution of \eqref{eq:codebook-opt} under a chosen constraint set.

\paragraph{We solve, not round.} With the grid fixed at $1/16$ ($m=4$) and the
endpoints pinned, the six interior levels are integers in $\{1,\dots,15\}$, so
\eqref{eq:codebook-opt} has only $\binom{15}{6}=5005$ feasible codebooks and is
solved \emph{exactly} by enumeration --- the reported codebook is the global
constrained optimum, not a rounding or a local (gradient/$k$-means) solution as
in learned-codebook methods. Crucially, we solve it on each model's \emph{own}
post-Hadamard activations (a short calibration pass), not on an assumed Gaussian.
The optimum $\mathcal{C}^\star=\tfrac{1}{16}\{0,2,4,6,8,10,13,16\}$ has $17\%$
lower distortion than GF4 rounded to the same grid
($\tfrac{1}{16}\{0,1,3,5,6,8,11,16\}$), which collapses adjacent levels at coarser
grids; and it is remarkably \emph{stable} across models. Solved independently on
each model's own activations over eight models and three families (OPT-125M–13B,
Mistral-7B, Qwen2.5-7B/14B), the shift-and-add optimum
$\tfrac{1}{16}\{0,2,4,6,8,10,12,16\}$ is \emph{identical for all eight}, and the
unconstrained optimum for seven (Qwen2.5-7B differs by a single grid step,
$13\!\to\!12$); the first six interior levels never vary. A single fixed-point
lattice codebook therefore suffices --- no per-layer or per-model lookup table,
unlike GPU-adaptive schemes such as GANQ~\cite{zhao2025ganq}. Notably, the
hardware constraint \emph{regularizes} the optimum: the popcount-constrained
codebook is perfectly stable where the unconstrained one wobbles by one level.

\paragraph{Results.} Table~\ref{tab:codebook-design} reports distortion, mean
perplexity gap, and decode-ROM area for the candidate codebooks. Over the
eight-model, three-family suite in the deployed W4A4 regime (NVFP4 weights $+$
codebook activations), the constrained optima match or slightly beat hand-tuned
GF4 on average --- mean $\Delta$PPL $-0.06$ (Q1.4) and $-0.09$
(popcount$\le2$), with the full per-model breakdown in
Table~\ref{tab:codebook-suite} --- while cutting decode-ROM area by $53\%$ and $61\%$
respectively; the shift-and-add codebook $\tfrac{1}{16}\{0,2,4,6,8,10,12,16\}$ is
both the cheapest to decode and the best on average. Naive rounding, by contrast,
is the worst grid codebook on every model (Figure~\ref{fig:codebook-suite}a). The
value is thus hardware at iso-accuracy: the constrained optimum matches the
distribution-aware codebook while being natively cheap to decode, and
Figure~\ref{fig:codebook-suite}b shows the solved codebook is the same lattice for
every model save a single one-step deviation.

\begin{figure}[t]
  \centering
  \includegraphics[width=\linewidth]{figures/gf4_codebook_suite.png}
  \caption{The constrained-optimal codebook over an eight-model, three-family
  suite (W4A4). \textbf{(a)} Perplexity gap vs.\ hand-tuned GF4 against model
  scale: the optimized codebooks track the GF4 baseline at every scale (dotted
  lines are means, both below zero), while naive rounding sits above it.
  \textbf{(b)} The codebook solved independently on each model's own activations
  is the same $k/16$ lattice for all eight models; the only variation is
  Qwen2.5-7B's penultimate level ($13\!\to\!12$, one grid step) --- so a single
  fixed-point codebook suffices, with no per-model lookup table.}
  \label{fig:codebook-suite}
\end{figure}

\paragraph{Measured GPU decode speedup.} The lattice constraint pays off on
today's GPUs, not only on the ASIC. We add a decode path to the fused
W4A16 GEMV of Section~\ref{sec:sw-support} that decodes the optimal codebook by
shift-and-add ($k=2\,\mathrm{idx}$, so $\ell=k/16$), touching no table, and
benchmark it against GF4's arbitrary codebook decoded from a lookup table.
Table~\ref{tab:gpu-speedup} reports the result on an RTX 4080: two effects
compose --- the $3.6\times$ weight-traffic reduction any 4-bit format enjoys over
dense fp16 ($1.5$--$1.7\times$), and a further $1.6$--$1.9\times$ from the
table-free decode --- for a net $2.6$--$3.1\times$ over dense fp16 in the
memory-bound (cold-L2) regime. The decode win is mechanistic: an arbitrary
codebook must be read from a table, and a constant-memory lookup with
\emph{divergent} codes across a warp serializes (constant memory is
broadcast-optimized), a cost that persists even when the kernel is otherwise
memory-bound; the lattice codebook is pure arithmetic and pays none of it. The
comparison is against a constant-memory table --- a register- or shared-resident
LUT would narrow the gap --- but the lattice is fundamentally table-free, so no
lookup implementation can be faster. This is the deployment counterpart to the
ASIC decode savings: the same property that shrinks the ROM (Table~\ref{tab:node-scaling})
removes the lookup on the GPU.

\paragraph{Relation to prior codebooks.} Distribution-aware 4-bit formats ---
NF4~\cite{dettmers2023qlora} and the optimized AF4~\cite{yoshida2023af4} and
BOF4~\cite{blumenberg2026bof4} --- optimize the codebook for the weight
distribution but place levels on the real line, giving arbitrary lookup constants
(as does the learned any4~\cite{elhoushi2025any4}). Hardware-efficient formats ---
additive powers-of-two (APoT)~\cite{li2020apot} and bit-vector sums
(SBVR)~\cite{bang2025sbvr} --- constrain levels to shift-friendly values but are
not derived from the operand distribution. GPU-adaptive schemes such as
GANQ~\cite{zhao2025ganq} optimize an arbitrary per-layer LUT. Our formulation
\eqref{eq:codebook-opt} is distribution-aware \emph{and} hardware-constrained,
and (per the universality above) a single lattice codebook rather than a per-layer
table; we found no prior 4-bit codebook that is a Gaussian-optimal solution on a
fixed dyadic lattice. That is what keeps the format portable across GPUs without
Blackwell-specific numeric support, while remaining a provable optimum rather than
a hand-tuned table.
